\subsection{正规化子}

设 \(G\)\(g\)\(G\)\(a\)\(g\)\(N_G(A)\) 是李群，g 是其李代数，A 是 G 的子集，a 是 g 的子集。在本节中，\(g \in G\) 表示满足 \(gAg^{-1} = A\) 的 \(G\) 的集合。它是 G 的子群；当 A 闭时它是闭子群。\(A\)\(n_g(a)\) 表示满足 \(x \in g\) 的 \([x, a] \subset a\) 的集合（参见第一章，§ 1，第 4 号）。它是 g 的子代数；当 a 闭时它是闭子代数。\(g\)\(a\)\(N_G(a)\) 表示满足 \(g \in G\) 的 \(gag^{-1} = a\) 的集合。

命题 10. 设 \(G\)\(g\)\(a\)\(g\)\(N_G(a)\) 是有限维李群，g 是其李代数，a 是 g 的向量子空间。则 \(G\) 是 G 的李子群，其李代数为 \(n_g(a)\)。

这由 § 3，命题 44 以及命题 39 的推论 1 得到。

命题 11. 设 G 是有限维实或复李群，g 是其李代数，A 是 G 的积分子群，且 \(\mathfrak{a} = \mathrm{L}(\mathbf{A})\)。则 \(\mathrm{N}_{\mathbf{G}}(\mathbf{A}) = \mathrm{N}_{\mathbf{G}}(\mathfrak{a})\)，且 \(\mathrm{N}_{\mathbf{G}}(\mathbf{A})\) 是包含 \(\overline{\mathbf{A}}\) 的 G 的李子群，其李代数为 \(\mathfrak{n}_{\mathfrak{g}}(\mathfrak{a})\)。

等式 \(\mathrm{N}_{\mathsf{G}}(\mathsf{A}) = \mathrm{N}_{\mathsf{G}}(\mathsf{a})\) 由 § 6，第 2 号，命题 3 的推论 2 得到。由命题 10，\(\mathrm{N}_{\mathsf{G}}(\mathsf{A})\) 是 G 的李子群，其李代数为 \(\mathfrak{n}_{\mathsf{g}}(\mathsf{a})\)。因此 \(\mathrm{N}_{\mathsf{G}}(\mathsf{A})\) 是闭的。由于 \(\mathrm{N}_{\mathsf{G}}(\mathsf{A}) \supset \mathsf{A}\)，有 \(\mathrm{N}_{\mathsf{G}}(\mathsf{A}) \supset \overline{\mathsf{A}}\)。

推论。若 \(\mathfrak{a} = \mathfrak{n}_{\mathfrak{g}}(\mathfrak{a})\)，则 A 是 G 的李子群，并且是 \(\mathrm{N}_{\mathbf{G}}(\mathbf{A})\) 的单位分支。

该单位分支是李子群，其李代数为 \(n_{g}(a)\)（命题 11），因而根据 §6，第 2 号，定理 2 (i)，它等于 A。
% semantic-fragments: legacy-input-seam
\relax{}
